\begin{center}\itshape
Pick an integer at random. Is it abundant? How close is $\varphi(n)$ to $n$?\\
The primes dividing $n$ behave like independent coins,\\
and an additive function is the sum of what they show.
\end{center}

In 1928 Isaac Schoenberg, a young Romanian mathematician then working in Germany, proved a curious fact about Euler's function: the proportion of integers $n\le x$ with $\frac{\varphi(n)}n\le z$ tends, for every $z$, to a limit $F(z)$, and $F$ is continuous. Five years later Harold Davenport proved the same kind of statement for $\frac{\sigma(n)}n$, the sum of the divisors of $n$ divided by $n$. As a consequence the \emph{abundant numbers}, those with $\sigma(n)\ge2n$, have a natural density, a question that went back to the Greeks' fascination with perfect numbers. In the same year Felix Behrend showed that the density lies between $0.241$ and $0.314$. It is now known to be $0.24761\ldots$, which means that roughly one integer in four is abundant.

These were hard-won theorems, each proved with its own machinery. In 1939 Paul Erd\H{o}s and Aurel Wintner, in a nine-page paper in the \emph{American Journal of Mathematics}, explained them all at once. Write $\log\frac{\varphi(n)}n=\sum_{p\mid n}\log(1-\frac1p)$. This is an \emph{additive} function: a sum over the primes dividing $n$. Erd\H{o}s and Wintner gave a criterion for an additive function to have a limiting distribution, and the criterion is the three-series theorem that Kolmogorov had proved in 1928 for sums of independent random variables. The title of their paper says why: \emph{Additive arithmetical functions and statistical independence}.

The independence is a theorem and not just a heuristic. Twenty years later Jonas Kubilius in Vilnius made it precise: the exponents of the small primes in a random integer $n\le x$ are, up to a tiny error in total variation, independent geometric random variables. Kubilius became rector of Vilnius University in 1958 and led it for more than thirty years, until 1990, while building the Lithuanian school of probability and number theory. His monograph, translated by the American Mathematical Society in 1964 as \emph{Probabilistic Methods in the Theory of Numbers}, gave the subject its name.

This chapter builds the model and draws the consequences: the Tur\'an--Kubilius inequality for every additive function, the Kubilius model in total variation, the Erd\H{o}s--Wintner theorem, and two theorems of probability, of L\'evy and of Jessen and Wintner, which decide whether the limit law is continuous. Euler's function and the abundant numbers come out as corollaries.

\begin{reminder}[from Volumes~4, 6, 11 and~14]
\emph{Tur\'an--Kubilius} (Volume~6, Theorem~9.1). There is an absolute $C_0$ such that for every real function $f$ on the primes and $x\ge3$, the strongly additive $f(n)=\sum_{p\mid n}f(p)$ satisfies $\sum_{n\le x}(f(n)-A_0(x))^2\le C_0xB_0(x)^2$, where $A_0(x)=\sum_{p\le x}\frac{f(p)}p$ and $B_0(x)^2=\sum_{p\le x}\frac{f(p)^2}p$.
\emph{Rough numbers} (Volume~14, Corollary~3.5). For $x\ge4$ and $2\le z\le\sqrt x$, the number $\Phi(x,z)$ of $n\le x$ with no prime factor $\le z$ is $x\prod_{p\le z}(1-\frac1p)(1+O(e^{-u/2}))+O(\sqrt x\log x)$, with $u=\frac{\log x}{\log z}$.
\emph{Three-series theorem} (Volume~11, Theorem~2.5). For independent $X_p$ and $c>0$, $\sum X_p$ converges almost surely if and only if $\sum P(|X_p|>c)$, $\sum\mathbb E[X_p\mathbf 1_{|X_p|\le c}]$ and $\sum\operatorname{Var}[X_p\mathbf 1_{|X_p|\le c}]$ converge.
\emph{Kolmogorov's zero--one law} (Volume~11, Theorem~1.4) and the \emph{Borel--Cantelli lemmas} (Volume~4, Chapter~4).
\emph{Mertens} (Volume~6, Theorems~1.3--1.6). $\sum_{p\le y}\frac{\log p}p=\log y+O(1)$, $\sum_{p\le y}\frac1p=\log\log y+O(1)$ and $\prod_{p\le y}(1-\frac1p)\asymp\frac1{\log y}$.
\end{reminder}

\section{Additive functions}

\begin{definition}\label{def:additive}
A function $f:\N_{\ge1}\to\C$ is \emph{additive} if $f(mn)=f(m)+f(n)$ whenever $\gcd(m,n)=1$, so that $f(n)=\sum_{p^k\parallel n}f(p^k)$, where $p^k\parallel n$ means that $p^k$ divides $n$ and $p^{k+1}$ does not. It is \emph{strongly additive} if moreover $f(p^k)=f(p)$ for all $k\ge1$, and \emph{completely additive} if $f(mn)=f(m)+f(n)$ for all $m,n$.
\end{definition}

The number $\omega(n)$ of distinct prime factors is strongly additive, $\Omega(n)$ (with multiplicity) and $\log n$ are completely additive, and $\log\frac{\varphi(n)}n=\sum_{p\mid n}\log(1-\frac1p)$ is strongly additive. If $g$ is multiplicative and positive, $\log g$ is additive; so $\log\frac{\sigma(n)}n$ is additive, with $\log\frac{\sigma(p^k)}{p^k}=\log\frac{1-p^{-k-1}}{1-p^{-1}}$.

The object of study is the \emph{distribution} of $f$: for $n$ chosen uniformly in $\{1,\dots,\lfloor x\rfloor\}$, the law of the random variable $f(n)$. We write $P_x$ and $\mathbb E_x$ for probabilities and expectations under this uniform choice, so that $P_x(f(n)\le z)=\frac1{\lfloor x\rfloor}\#\{n\le x:f(n)\le z\}$. We say that $f$ has a \emph{limiting distribution} if $f(n)$ converges in law as $x\to\infty$, that is, if $P_x(f(n)\le z)$ tends to $F(z)$ at every continuity point $z$ of some distribution function $F$.

\begin{think}
Which of $\omega(n)$, $\log\frac{\varphi(n)}n$ and $\log\frac{\sigma(n)}n$ would you expect to have a limiting distribution without any renormalisation? Think of the size of the contribution of each prime, and of how often that prime divides $n$.
\end{think}

The heuristic is simple. A prime $p$ divides a random integer with probability $\frac1p$, and does so ``independently'' of the other primes. For $\omega$ each prime contributes $1$ with probability $\frac1p$, and $\sum\frac1p$ diverges: $\omega(n)$ grows like $\log\log n$ and must be centred and scaled, which is the Erd\H{o}s--Kac theorem of Volume~6 and of this chapter. For $\log\frac{\varphi(n)}n$ the prime $p$ contributes about $-\frac1p$ with probability $\frac1p$, and $\sum\frac1{p^2}$ converges: the sum should converge without any normalisation. The rest of the chapter makes this precise.

\section{The Tur\'an--Kubilius inequality}

Volume~6 proved the inequality for strongly additive functions. Every additive function differs from a strongly additive one only at the prime powers $p^k$ with $k\ge2$, which are rare, and that is the whole proof of the general case.

\begin{theorem}[Tur\'an--Kubilius]\label{thm:tkgeneral}
There is an absolute constant $C_1$ such that for every additive $f:\N\to\C$ and every $x\ge3$,
\[
\sum_{n\le x}|f(n)-A(x)|^2\le C_1xB(x)^2,\qquad A(x)=\sum_{p^k\le x}\frac{f(p^k)}{p^k}\Bigl(1-\frac1p\Bigr),\qq B(x)^2=\sum_{p^k\le x}\frac{|f(p^k)|^2}{p^k}.
\]
One can take $C_1=3C_0+36$.
\end{theorem}

The centring $A(x)$ is the expectation of the \emph{model}: if $G_p$ is a geometric variable with $P(G_p=k)=(1-\frac1p)p^{-k}$, then $\mathbb E f(p^{G_p})=\sum_{k\ge1}f(p^k)(1-\frac1p)p^{-k}$, with $f(1)=0$.

\begin{proof}
Since $|w|^2=(\operatorname{Re}w)^2+(\operatorname{Im}w)^2$ and the real and imaginary parts of $f$ are additive with $A$ and $B^2$ splitting accordingly ($B^2=B_{\operatorname{Re}}^2+B_{\operatorname{Im}}^2$), it suffices to treat real $f$ with the constant $3C_0+36$.

For $n\le x$ write $f(n)=g(n)+h(n)$ with
\[
g(n)=\sum_{p\mid n}f(p),\qquad h(n)=\sum_{\substack{p^k\parallel n\\k\ge2}}c(p^k),\qquad c(p^k)=f(p^k)-f(p).
\]
The function $g$ is strongly additive, and Volume~6, Theorem~9.1, gives $\sum_{n\le x}(g(n)-A_0)^2\le C_0xB_0^2\le C_0xB^2$, with $A_0=\sum_{p\le x}\frac{f(p)}p$.

\emph{The function $h$.} Expanding the square, the diagonal terms are $\sum_{p^k}c(p^k)^2\#\{n\le x:p^k\parallel n\}\le x\sum_{p^k}\frac{c(p^k)^2}{p^k}$. In the off-diagonal terms $p^k\parallel n$ and $q^l\parallel n$ with $p^k\ne q^l$ force $p\ne q$, so $p^kq^l\mid n$ and there are at most $\frac x{p^kq^l}$ such $n$. Hence, with all sums over $p^k\le x$ with $k\ge2$, and by Cauchy--Schwarz,
\[
\sum_{n\le x}h(n)^2\le x\sum\frac{c(p^k)^2}{p^k}+x\Bigl(\sum\frac{|c(p^k)|}{p^k}\Bigr)^2\le x\,H\Bigl(1+\sum_p\sum_{k\ge2}\frac1{p^k}\Bigr)\le2xH,
\]
where $H=\sum\frac{c(p^k)^2}{p^k}$, because $\sum_p\sum_{k\ge2}p^{-k}=\sum_p\frac1{p(p-1)}<\sum_{m\ge2}\frac1{m(m-1)}=1$. Now $c(p^k)^2\le2f(p^k)^2+2f(p)^2$, and $\sum_{k\ge2}p^{-k}=\frac1{p(p-1)}\le\frac1p$, so
\[
H\le2\sum_{\substack{p^k\le x\\k\ge2}}\frac{f(p^k)^2}{p^k}+2\sum_{p\le x}\frac{f(p)^2}p\le4B^2,\qquad\text{and}\qquad\sum_{n\le x}h(n)^2\le8xB^2.
\]

\emph{The centring.} Since $\frac{f(p)}p-\frac{f(p)}p(1-\frac1p)=\frac{f(p)}{p^2}$,
\[
A_0-A=\sum_{p\le x}\frac{f(p)}{p^2}-\sum_{\substack{p^k\le x\\k\ge2}}\frac{f(p^k)}{p^k}\Bigl(1-\frac1p\Bigr),
\]
and by Cauchy--Schwarz each sum is at most $B$ in absolute value: $|\sum_p\frac{f(p)}{p^2}|\le(\sum_p\frac{f(p)^2}p)^{1/2}(\sum_p\frac1{p^3})^{1/2}\le B$ and $|\sum\frac{f(p^k)}{p^k}|\le B(\sum_p\sum_{k\ge2}p^{-k})^{1/2}\le B$. So $|A_0-A|\le2B$.

Finally $f(n)-A=(g(n)-A_0)+h(n)+(A_0-A)$ and $(a+b+c)^2\le3(a^2+b^2+c^2)$ give $\sum_{n\le x}(f(n)-A)^2\le3(C_0+8+4)xB^2$.
\end{proof}

In the language of probability: $\operatorname{Var}_xf(n)\le\mathbb E_x|f(n)-A(x)|^2\le C_1B(x)^2(1+O(\frac1x))$, and $B(x)^2$ is, up to bounded factors, the variance of the model $\sum_{p\le x}f(p^{G_p})$. The dependence between the primes costs at most a constant factor. Chebyshev's inequality turns this into a law of large numbers: if $B(x)\to\infty$, then $\frac{f(n)-A(x)}{B(x)\,\psi(x)}\to0$ in probability for every $\psi(x)\to\infty$. For $\omega$ this is the theorem of Hardy and Ramanujan (Volume~6, Theorem~9.2).

\section{The Kubilius model}

The inequality compares only two moments. Kubilius's insight was that much more is true: for the primes up to a small power of $x$, the whole joint law of the exponents is close to the independent model, in the strongest sense available.

For $y\ge2$ and $n\ge1$ let $v_y(n)=(\nu_p(n))_{p\le y}$ be the vector of exponents of the primes $p\le y$ in $n$. Let $(G_p)_p$ be independent with $P(G_p=k)=(1-\frac1p)p^{-k}$ for $k\ge0$, and $G^{(y)}=(G_p)_{p\le y}$. Recall that the \emph{total variation distance} between two laws $\mu,\nu$ on a countable set is $d_{TV}(\mu,\nu)=\sup_E|\mu(E)-\nu(E)|=\frac12\sum_a|\mu(a)-\nu(a)|$.

\begin{theorem}[Kubilius]\label{thm:kubilius}
There is an absolute constant $C_2$ such that for $x\ge16$ and $2\le y\le x^{1/4}$, with $u=\frac{\log x}{\log y}$,
\[
d_{TV}\bigl(\operatorname{law}_x v_y(n),\ \operatorname{law}G^{(y)}\bigr)\le C_2\bigl(e^{-u/4}+x^{-1/4}\log x\bigr).
\]
\end{theorem}

\begin{proof}
A vector $\mathbf k=(k_p)_{p\le y}$ with finitely many non-zero entries corresponds to a $y$-\emph{smooth} integer $m=\prod_pp^{k_p}$, one all of whose prime factors are $\le y$. Write $V=\prod_{p\le y}(1-\frac1p)$, so $P(G^{(y)}=\mathbf k)=\frac Vm$. On the other side, $v_y(n)=\mathbf k$ means $n=mn'$ with $n'$ free of prime factors $\le y$, so
\[
P_x(v_y(n)=\mathbf k)=\frac{\Phi(x/m,y)}{\lfloor x\rfloor}.
\]
\emph{Small $m$.} Let $m\le\sqrt x$. Then $X=\frac xm\ge\sqrt x\ge y^2$, and Volume~14, Corollary~3.5, applies with $u_X=\frac{\log X}{\log y}\ge\frac u2$:
\[
\Bigl|\frac{\Phi(x/m,y)}x-\frac Vm\Bigr|\ll\frac Vme^{-u/4}+\frac{\sqrt{x/m}\log x}x.
\]
Summing over $y$-smooth $m\le\sqrt x$, and using $\sum_{m\ y\text{-smooth}}\frac1m=\prod_{p\le y}(1-\frac1p)^{-1}=\frac1V$ and $\sum_{m\le\sqrt x}m^{-1/2}\le2x^{1/4}$,
\[
\Sigma_1:=\sum_{m\le\sqrt x}\Bigl|\frac{\Phi(x/m,y)}x-\frac Vm\Bigr|\ll e^{-u/4}+x^{-1/4}\log x.
\]
\emph{Large $m$, model side (Rankin's trick).} Let $\delta=\frac1{\log y}$. For $m>\sqrt x$, $1\le(\frac m{\sqrt x})^\delta$, so
\[
\Sigma_2:=\sum_{\substack{m>\sqrt x\\y\text{-smooth}}}\frac Vm\le Vx^{-\delta/2}\sum_{y\text{-smooth}}\frac1{m^{1-\delta}}=Ve^{-u/2}\prod_{p\le y}\Bigl(1-\frac1{p^{1-\delta}}\Bigr)^{-1}.
\]
For $p\le y$, $p^\delta=e^{\log p/\log y}\le1+(e-1)\frac{\log p}{\log y}$ by convexity, so $\sum_{p\le y}p^{\delta-1}\le\sum_{p\le y}\frac1p+\frac{e-1}{\log y}\sum_{p\le y}\frac{\log p}p=\log\log y+O(1)$, and $\sum_pp^{2\delta-2}\le\sum_pe^2p^{-2}=O(1)$. With $\log(1-t)^{-1}\le t+t^2$ for $0\le t\le\frac12$ (and the factor at $p=2$ bounded), the product is $\ll\log y$, and $V\ll\frac1{\log y}$. Hence $\Sigma_2\ll e^{-u/2}$.

\emph{Large $m$, integer side.} The probabilities $P_x(v_y(n)=\mathbf k)$ sum to $1$, and so do the $\frac Vm$. Since $\lfloor x\rfloor\le x$,
\[
\sum_{m>\sqrt x}\frac{\Phi(x/m,y)}{\lfloor x\rfloor}=1-\sum_{m\le\sqrt x}\frac{\Phi(x/m,y)}{\lfloor x\rfloor}\le1-\sum_{m\le\sqrt x}\frac{\Phi(x/m,y)}x\le1-\sum_{m\le\sqrt x}\frac Vm+\Sigma_1=\Sigma_2+\Sigma_1.
\]
\emph{Conclusion.} Since $\sum_{m\le\sqrt x}\Phi(x/m,y)\le x$, replacing $\frac1x$ by $\frac1{\lfloor x\rfloor}$ in $\Sigma_1$ costs at most $\frac x{\lfloor x\rfloor}-1\le\frac2x$. So
\[
2d_{TV}=\sum_{\mathbf k}\bigl|P_x(v_y(n)=\mathbf k)-P(G^{(y)}=\mathbf k)\bigr|\le\Bigl(\Sigma_1+\frac2x\Bigr)+(\Sigma_1+\Sigma_2)+\Sigma_2\ll e^{-u/4}+x^{-1/4}\log x.
\]
\end{proof}

\begin{think}
With $y=x^{1/u}$, the theorem is useful only when $u$ is large. What goes wrong for $y=\sqrt x$? (Count the $n\le x$ that are prime, and compare with the model's probability that $G_p=0$ for every $p\le\sqrt x$.)
\end{think}

For $y=\sqrt x$ the integer side has $P_x(v_y(n)=\mathbf0)=\frac{1+\pi(x)-\pi(\sqrt x)}{\lfloor x\rfloor}\sim\frac1{\log x}$, while the model gives $V\sim\frac{2e^{-\gamma}}{\log x}$: the two differ by the factor $\frac{e^\gamma}2\approx0.89$ already met in Volume~14, Figure~3.1. A large prime factor uses up most of the room below $x$, and that constraint is invisible to independent coins. The model is exact for the primes up to $x^{o(1)}$, and only for those.

Since $f(n)$ restricted to the small primes is a function of $v_y(n)$, the theorem transfers at once to every additive function. Let
\[
f_y(n)=\sum_{\substack{p^k\parallel n\\p\le y}}f(p^k),\qquad Z_y=\sum_{p\le y}f(p^{G_p})\quad(f(1)=0).
\]

\begin{corollary}\label{cor:kubiliusf}
For every additive $f$, $x\ge16$ and $2\le y\le x^{1/4}$, and every Borel set $E\subseteq\C$, $|P_x(f_y(n)\in E)-P(Z_y\in E)|\le C_2(e^{-u/4}+x^{-1/4}\log x)$.
\end{corollary}
\begin{proof}
$f_y(n)=\Lambda(v_y(n))$ and $Z_y=\Lambda(G^{(y)})$ for the same map $\Lambda(\mathbf k)=\sum_pf(p^{k_p})$, and total variation distance can only decrease under a map: $\{f_y(n)\in E\}=\{v_y(n)\in\Lambda^{-1}(E)\}$.
\end{proof}

\section{The Erd\H{o}s--Wintner theorem}

For a real additive $f$ and a prime $p$ write $f^*(p)=f(p)$ if $|f(p)|\le1$ and $f^*(p)=0$ otherwise.

\begin{theorem}[Erd\H{o}s and Wintner, 1939]\label{thm:erdoswintner}
Let $f$ be a real additive function such that the three series
\[
\sum_{|f(p)|>1}\frac1p,\qquad\sum_p\frac{f^*(p)}p,\qquad\sum_p\frac{f^*(p)^2}p
\]
converge. Then $f$ has a limiting distribution: $f(n)$, for $n$ uniform in $[1,x]$, converges in law as $x\to\infty$ to
\[
Z=\sum_pf(p^{G_p}),
\]
a series of independent variables that converges almost surely. Its characteristic function is
\[
\mathbb Ee^{itZ}=\prod_p\Bigl(1-\frac1p\Bigr)\Bigl(1+\sum_{k\ge1}\frac{e^{itf(p^k)}}{p^k}\Bigr).
\]
\end{theorem}

\begin{proof}
\emph{The model converges.} Let $X_p=f(p^{G_p})$ and apply the three-series theorem with $c=1$. Since $P(G_p\ge2)=p^{-2}$,
\[
\begin{gathered}P(|X_p|>1)\le\frac1{p^2}+\frac{\mathbf 1_{|f(p)|>1}}p,\\\Bigl|\mathbb E[X_p\mathbf 1_{|X_p|\le1}]-f^*(p)\frac{p-1}{p^2}\Bigr|\le\frac1{p^2},\qquad\mathbb E[X_p^2\mathbf 1_{|X_p|\le1}]\le\frac{f^*(p)^2}p+\frac1{p^2}.\end{gathered}
\]
The three hypotheses, with $\sum_p\frac1{p^2}<\infty$ and $|\frac{f^*(p)}{p^2}|\le\frac1{p^2}$, make the three series of Volume~11, Theorem~2.5, converge (the variance is at most the second moment). So $Z_y\to Z$ almost surely as $y\to\infty$, hence in law. Since $Z$ is the almost sure limit of independent sums, $\mathbb Ee^{itZ}=\lim_y\prod_{p\le y}\mathbb Ee^{itX_p}$ by dominated convergence, which is the product in the statement.

\emph{The small primes.} Let $y=x^{1/\log\log x}$, so $u=\log\log x\to\infty$ and $y\to\infty$. By Corollary~\ref{cor:kubiliusf} the laws of $f_y(n)$ and $Z_y$ are $o(1)$ apart in total variation. Since $Z_y\to Z$ in law, so does $f_y(n)$: for every bounded continuous $\phi$, $|\mathbb E_x\phi(f_y(n))-\mathbb E\phi(Z_y)|\le2\|\phi\|_\infty d_{TV}\to0$.

\emph{The large primes.} Let $R(n)=f(n)-f_y(n)=\sum_{p^k\parallel n,\,p>y}f(p^k)$. We show that $R(n)\to0$ in probability; Slutsky's lemma (Volume~4, Exercise~6.1) then gives $f(n)=f_y(n)+R(n)\to Z$ in law. Let $\tilde f$ be the strongly additive function with $\tilde f(p)=f^*(p)$ for $y<p\le x$ and $\tilde f(p)=0$ otherwise. Then $R(n)=\tilde f(n)$ unless $p^2\mid n$ for some $p>y$, or $p\mid n$ for some $p>y$ with $|f(p)|>1$. So
\[
P_x(R(n)\ne\tilde f(n))\le\sum_{p>y}\frac1{p^2}+\sum_{\substack{p>y\\|f(p)|>1}}\frac1p+O\Bigl(\frac{\pi(x)}x\Bigr)\longrightarrow0,
\]
as tails of convergent series (the last term accounts for the integer parts $\lfloor x/p\rfloor$ divided by $\lfloor x\rfloor$). By Volume~6, Theorem~9.1, applied to $\tilde f$, with $\tilde A=\sum_{y<p\le x}\frac{f^*(p)}p$ and $\tilde B^2=\sum_{y<p\le x}\frac{f^*(p)^2}p$,
\[
P_x\bigl(|\tilde f(n)-\tilde A|>\varepsilon\bigr)\le\frac{C_0\tilde B^2}{\varepsilon^2}\Bigl(1+O\Bigl(\frac1x\Bigr)\Bigr)\longrightarrow0,
\]
and $\tilde A\to0$, $\tilde B\to0$ as tails of the second and third series. So $P_x(|R(n)|>2\varepsilon)\to0$ for every $\varepsilon>0$.
\end{proof}

\begin{remark}[the converse]
Erd\H{o}s and Wintner also proved the converse: if $f$ has a limiting distribution, the three series converge. The proof needs a lower bound for characteristic functions of additive functions that we do not develop; see Tenenbaum's book, Chapter~III.4, or Elliott's \emph{Probabilistic Number Theory~I}, Chapter~5. Notice the asymmetry: the prime powers $p^k$ with $k\ge2$ play no role in the criterion, because $\sum_p\sum_{k\ge2}p^{-k}<\infty$.
\end{remark}

For $f(n)=\log\frac{\varphi(n)}n$ we have $f(p)=\log(1-\frac1p)\in[-\log2,0[$, so $f^*(p)=f(p)=-\frac1p+O(\frac1{p^2})$, and the three series converge (the first is empty). The same holds for $\log\frac{\sigma(n)}n$, with $f(p)=\log(1+\frac1p)$. So both have limiting distributions, and so do $\frac{\varphi(n)}n$ and $\frac{\sigma(n)}n$, by continuity of $\exp$. Figure~\ref{fig:phi} shows the first.

\begin{figure}[ht]
\centering\fig{ch01-phi}
\caption{The proportion of $n\le10^6$ with $\frac{\varphi(n)}n\le z$, against $z$. The limit law exists by Theorem~\ref{thm:erdoswintner} and is continuous by Corollary~\ref{cor:phicont}; the graph has no jumps, but its steep climbs at $z=\frac12$ (the even numbers not divisible by~$3$), $z=\frac13$ (multiples of $6$ prime to~$5$) and their neighbours show how singular it is. The mean is $\frac6{\pi^2}=0.6079$ (Exercise~\ref{ex:meanphi}).}
\label{fig:phi}
\end{figure}

\section{Continuity and purity of the limit law}

A sum of independent discrete variables is discrete when it has finitely many terms. With infinitely many it may be continuous: $\sum_k\varepsilon_k2^{-k}$, with independent fair $\varepsilon_k\in\{0,1\}$, is uniform on $[0,1]$. Paul L\'evy found in 1931 exactly when this happens.

\begin{think}
The limit law of $\frac{\varphi(n)}n$ is the law of $\prod_p(1-\frac1p)^{B_p}$, with independent $B_p$ equal to $1$ with probability about $\frac1p$. Does it have an atom, a value taken with positive probability? Compare with $\omega$-like functions where $f(p)=0$ for most primes.
\end{think}

The key is a lemma from the geometry of sums of signs.

\begin{lemma}[Littlewood--Offord, in Erd\H{o}s's form, 1945]\label{lem:lo}
Let $a_1,\dots,a_m$ be non-zero real numbers and $\varepsilon_1,\dots,\varepsilon_m$ independent with $P(\varepsilon_i=1)=P(\varepsilon_i=-1)=\frac12$. Then for every real $c$,
\[
P\Bigl(\sum_{i=1}^m\varepsilon_ia_i=c\Bigr)\le\binom m{\lfloor m/2\rfloor}2^{-m}\le\frac1{\sqrt m}.
\]
\end{lemma}
\begin{proof}
Replacing $(\varepsilon_i,a_i)$ by $(-\varepsilon_i,-a_i)$ does not change the law, so we may assume every $a_i>0$. A choice of signs is a subset $I=\{i:\varepsilon_i=1\}$ of $\{1,\dots,m\}$, and $\sum\varepsilon_ia_i=2\sum_{i\in I}a_i-\sum_ia_i$. If $I\subsetneq J$, the sum for $J$ exceeds the sum for $I$ by $2\sum_{J\setminus I}a_i>0$. So the sets $I$ giving the value $c$ form an \emph{antichain}: none contains another.

\emph{Sperner's bound.} An antichain $\mathcal F$ of subsets of an $m$-set has at most $\binom m{\lfloor m/2\rfloor}$ members. Indeed, consider the $m!$ maximal chains $\emptyset\subset\{\sigma(1)\}\subset\{\sigma(1),\sigma(2)\}\subset\cdots$, one for each permutation $\sigma$. A set $I$ of size $j$ lies on exactly $j!(m-j)!$ of them, and each chain meets the antichain at most once. Hence $\sum_{I\in\mathcal F}|I|!(m-|I|)!\le m!$, that is $\sum_{I\in\mathcal F}\binom m{|I|}^{-1}\le1$, and each $\binom m{|I|}\le\binom m{\lfloor m/2\rfloor}$.

\emph{The binomial bound.} For $m=2k$, $\bigl(\binom{2k}k4^{-k}\bigr)^2=\prod_{j=1}^k\bigl(\frac{2j-1}{2j}\bigr)^2\le\prod_{j=1}^k\frac{2j-1}{2j}\cdot\frac{2j}{2j+1}=\frac1{2k+1}$, so $\binom{2k}k4^{-k}\le\frac1{\sqrt{2k+1}}\le\frac1{\sqrt m}$. For $m=2k+1$, $\binom{2k+1}k2^{-2k-1}=\binom{2k}k4^{-k}\frac{2k+1}{2k+2}\le\frac1{\sqrt{2k+1}}=\frac1{\sqrt m}$.
\end{proof}

\begin{theorem}[L\'evy, 1931]\label{thm:levyatoms}
Let $X_1,X_2,\dots$ be independent discrete real random variables such that $S=\sum_jX_j$ converges almost surely, and let $q_j=\max_vP(X_j=v)$. Then the law of $S$ has an atom if and only if $\sum_j(1-q_j)<\infty$; otherwise it is continuous.
\end{theorem}

\begin{proof}
\emph{Convergent case.} Choose $v_j$ with $P(X_j=v_j)=q_j$. By the first Borel--Cantelli lemma, almost surely $X_j=v_j$ for all large $j$. Since $\sum X_j$ converges almost surely, $\sum_jv_j$ converges (its tail coincides with a tail of $\sum X_j$ on an event of positive probability). The events $E_N=\{X_j=v_j\ \forall j>N\}$ increase to an almost sure event, so $P(E_N)>0$ for some $N$. The finite sum $S_N=X_1+\dots+X_N$ is discrete, so $P(S_N=w)>0$ for some $w$. By independence,
\[
P\Bigl(S=w+\sum_{j>N}v_j\Bigr)\ge P(S_N=w)P(E_N)>0.
\]
\emph{Divergent case.} Let $S'$ be an independent copy of $S$, built from an independent copy $(X'_j)$. For every $a$,
\[
P(S=a)^2=P(S=a,S'=a)\le P(S-S'=0).
\]
Now $S-S'=\sum_jD_j$ with $D_j=X_j-X'_j$ independent and symmetric, and $P(D_j\ne0)=1-\sum_vP(X_j=v)^2\ge1-q_j$ since $\sum_vP(X_j=v)^2\le q_j\sum_vP(X_j=v)=q_j$. So $\sum_jP(D_j\ne0)=\infty$, and by the second Borel--Cantelli lemma infinitely many $D_j$ are non-zero, almost surely. Let $(\varepsilon_j)$ be independent fair signs, independent of everything else. By symmetry and independence, the sequence $(\varepsilon_jD_j)$ has the same law as $(D_j)$, so $P(\sum_jD_j=0)=P(\sum_j\varepsilon_jD_j=0)$.

Fix $m\ge1$ and $N$ so large that the event $F$ that at least $m$ of $D_1,\dots,D_N$ are non-zero has $P(F)\ge1-\frac1m$. Condition on all the $D_j$ and on the $\varepsilon_j$ with $D_j=0$ or $j>N$. On $F$, the sum $\sum_j\varepsilon_jD_j$ is then $c+\sum_{j\in J}\varepsilon_jD_j$ for a fixed $c$ and a set $J$ of at least $m$ indices with $D_j\ne0$, and the remaining signs are still independent and fair. Lemma~\ref{lem:lo} gives a conditional probability at most $\frac1{\sqrt m}$ that the sum is $0$. So $P(\sum D_j=0)\le\frac1{\sqrt m}+\frac1m$ for every $m$, hence $P(\sum D_j=0)=0$ and $S$ has no atom.
\end{proof}

\begin{corollary}\label{cor:ewcont}
Under the hypotheses of Theorem~\ref{thm:erdoswintner}, the limit law of $f$ is continuous if and only if $\sum_{f(p)\ne0}\frac1p=\infty$.
\end{corollary}
\begin{proof}
Let $X_p=f(p^{G_p})$. For $p\ge3$ every value other than $0$ is taken with probability at most $P(G_p\ge1)=\frac1p<1-\frac1p\le P(X_p=0)$, so $q_p=P(X_p=0)$ and $1-q_p=P(X_p\ne0)$. If $f(p)\ne0$ then $P(X_p\ne0)\ge P(G_p=1)=\frac{p-1}{p^2}\ge\frac1{2p}$; and always $P(X_p\ne0)\le\frac{\mathbf 1_{f(p)\ne0}}p+\frac1{p^2}$. So $\sum_p(1-q_p)$ and $\sum_{f(p)\ne0}\frac1p$ converge or diverge together, and Theorem~\ref{thm:levyatoms} applies.
\end{proof}

\begin{corollary}[Schoenberg 1928, Davenport 1933]\label{cor:phicont}
The functions $\frac{\varphi(n)}n$ and $\frac{\sigma(n)}n$ have continuous limiting distribution functions. In particular the set of abundant numbers, $\{n:\sigma(n)\ge2n\}$, has a natural density.
\end{corollary}
\begin{proof}
For both logarithms $f(p)\ne0$ for every $p$, and $\sum\frac1p=\infty$. Continuity of the limit law $F$ of $\log\frac{\sigma(n)}n$ at $\log2$ means that $P_x(\log\frac{\sigma(n)}n\ge\log2)\to1-F(\log2^-)=1-F(\log2)$.
\end{proof}

The argument shows more: the limit law of $\log\frac{\sigma(n)}n$ gives zero mass to every single point, so the density of $\{n:\sigma(n)\ge\alpha n\}$ exists for every $\alpha$ and depends continuously on $\alpha$. Figure~\ref{fig:sigma} shows the histogram up to $10^6$: $24.75\%$ of the integers are abundant there. The limit is $0.24761\ldots$, computed by Deléglise and then, to more digits, by Kobayashi and others, from the characteristic function of Theorem~\ref{thm:erdoswintner}.

\begin{figure}[ht]
\centering\fig{ch01-sigma}
\caption{Histogram of $\frac{\sigma(n)}n$ for $n\le10^6$ (density scale). The spike at $\frac{\sigma(n)}n=1^+$ comes from the primes and the integers with only large prime factors; abundant numbers are those right of the rust line. The largest value below $10^6$ is $4.509$, at $n=720720$.}
\label{fig:sigma}
\end{figure}

Continuity is not the end of the story. A probability law on $\R$ splits uniquely as $\mu=\mu_d+\mu_{sc}+\mu_{ac}$, a discrete part, a continuous part carried by a Lebesgue-null set, and an absolutely continuous part (Volume~7, Theorem~11.5; the singular part then splits into its atoms and the rest). Jessen and Wintner proved in 1935 that laws of convergent sums of independent discrete variables are always of one pure type.

\begin{theorem}[Jessen and Wintner, 1935]\label{thm:jessenwintner}
Under the hypotheses of Theorem~\ref{thm:levyatoms}, the law $\mu$ of $S$ is either purely discrete, or purely singular continuous, or absolutely continuous.
\end{theorem}
\begin{proof}
Let $G$ be the additive subgroup of $\R$ generated by all the values $v$ with $P(X_j=v)>0$, for all $j$; it is countable. For a Borel set $B$, the set $B+G=\bigcup_{g\in G}(B+g)$ is Borel. Almost surely every partial sum $S_N$ lies in $G$, so $S\in B+G$ if and only if $\sum_{j>N}X_j\in B+G$. This holds for every $N$, so up to a null set $\{S\in B+G\}$ is a tail event, and $\mu(B+G)\in\{0,1\}$ by Kolmogorov's zero--one law.

If $\mu_d\ne0$, let $D$ be the countable set of atoms; then $\mu(D+G)\ge\mu(D)>0$, so $\mu(D+G)=1$, and $D+G$ is countable: $\mu$ is purely discrete. If $\mu_d=0$ and $\mu_{sc}\ne0$, let $N$ be a Lebesgue-null set with $\mu_{sc}(N)=\mu_{sc}(\R)$. Then $N+G$ is null, a countable union of null sets, so $\mu_{ac}(N+G)=0$; and $\mu(N+G)\ge\mu_{sc}(N)>0$ forces $\mu(N+G)=1$, hence $\mu_{ac}(\R)=\mu_{ac}(\R)-\mu_{ac}(N+G)\le\mu(\R\setminus(N+G))=0$. So $\mu=\mu_{sc}$. Otherwise $\mu=\mu_{ac}$.
\end{proof}

\begin{remark}[which pure type?]
Erd\H{o}s showed in 1939 that the limit laws of $\frac{\varphi(n)}n$ and $\frac{\sigma(n)}n$ are purely singular: their distribution functions are continuous but have derivative $0$ almost everywhere. The steep climbs in Figure~\ref{fig:phi} are the visible trace of this. Deciding the type of the law of $\sum_p f(p^{G_p})$ in general is hard, and is related to the theory of Bernoulli convolutions; see Elliott, \emph{Probabilistic Number Theory~I}, Chapter~7.
\end{remark}

\begin{beauty}
Two unrelated-looking facts, that one integer in four is abundant and that $\frac{\varphi(n)}n$ has a continuous distribution, are instances of one probabilistic theorem about sums of independent coins. The arithmetic enters only through two counts: a prime $p$ divides a proportion $\frac1p$ of the integers, and the divisibilities by different small primes are independent up to an error that the sieve controls. Everything else is probability, and it is the probability of Kolmogorov, L\'evy and Sperner, a combinatorialist's lemma about subsets.
\end{beauty}
